% 将本文件中的说明替换为本次实验正文。章节可按课程要求增删。
\section{实验背景与任务说明}
说明实验背景、实验目的、题目要求及完成范围。

\section{实验成员分工}
个人实验可删除本节。
\begin{table}[H]
\centering
\caption{小组成员分工}
\begin{tabularx}{\textwidth}{p{3cm}X}
\toprule 成员 & 主要工作\\\midrule
姓名一 & 填写实际承担的任务。\\
姓名二 & 填写实际承担的任务。\\
\bottomrule
\end{tabularx}
\end{table}

\section{实验环境与工具}
\begin{table}[H]
\centering
\caption{实验环境}
\begin{tabularx}{\textwidth}{p{3cm}X}
\toprule 项目 & 配置或工具\\\midrule
操作系统 & 填写系统及版本。\\
硬件或虚拟机 & 填写实际配置。\\
编程语言与依赖 & 填写语言、编译器及依赖版本。\\
开发与调试工具 & 填写工具及版本。\\
\bottomrule
\end{tabularx}
\end{table}

\section{实验原理与方案设计}
\subsection{基本原理}
说明关键概念、算法或系统机制。公式示例（应替换为本次实验的实际公式）：
\begin{equation}
\bar{t}=\frac{1}{n}\sum_{i=1}^{n}t_i
\label{eq:mean}
\end{equation}

\subsection{整体流程}
以下为版式示例，按实际实验修改节点和连接关系。
\begin{figure}[H]
\centering
\begin{tikzpicture}[
  node distance=0.6cm,
  stage/.style={draw,rounded corners=3pt,text width=8cm,
    minimum height=0.8cm,align=center,fill=black!4},
  arr/.style={-{Stealth[length=2.2mm]},thick}]
\node[stage] (prepare) {环境准备与输入检查};
\node[stage,below=of prepare] (run) {执行实验步骤};
\node[stage,below=of run] (verify) {记录结果与核验};
\draw[arr] (prepare) -- (run);
\draw[arr] (run) -- (verify);
\end{tikzpicture}
\caption{实验流程示例}
\end{figure}

\section{实验步骤与关键实现}
\subsection{环境准备}
填写安装、配置、构建或启动步骤，以及实际使用的命令。
\subsection{关键代码}
以下代码仅演示代码块样式，不代表本次实验的实现。代码中建议使用英文注释；中文解释放在代码块外。
\begin{lstlisting}[language=Python,caption={代码排版示例}]
def mean_time(samples):
    if not samples:
        raise ValueError("samples must not be empty")
    return sum(samples) / len(samples)
\end{lstlisting}

\section{实验结果与分析}
\subsection{运行结果}
先描述实际现象，再插入对应截图并解释其能支持的结论。
\reportfigure{figures/result.png}{实验运行结果（待替换为真实截图）}

\subsection{测试与数据分析}
下表中的“待填写”为占位内容。填写实际数据、单位、测试条件和重复次数后再作分析。
\begin{table}[H]
\centering
\caption{测试结果记录}
\label{tab:results}
\begin{tabularx}{\textwidth}{p{2.5cm}p{2.6cm}X}
\toprule 测试项 & 实际结果 & 现象与分析\\\midrule
正常流程 & 待填写 & 记录实际执行现象。\\
边界或异常情况 & 待填写 & 记录输入条件、结果与原因。\\
性能指标 & 待填写 & 写明单位、环境与统计方式。\\
\bottomrule
\end{tabularx}
\end{table}
交叉引用示例：表\ref{tab:results}记录实际结果，式\eqref{eq:mean}演示公式引用。

\section{问题讨论与改进方向}
记录实际遇到的问题、定位依据、解决方法，以及尚未解决的限制。

\section{实验总结}
总结完成了哪些目标、哪些结果有证据支持、获得了什么认识，以及可继续改进的方向。
